% Split from 5_nina.tex by cut-sheet v2/v3 — content below is the source scene, header adjusted
\chapter{The Almost Integer}
\begin{minipage}[t]{\textwidth}
\lettrine{N}{ina} Ashtekar leaned back in her chair, the sunlight of a Rajastani hotel streaming through carved lattice windows.  

\begin{figure}[H]
\includegraphics[width=\linewidth]{Illustrations/vign-ninaAshketar.png}
\caption{Nina setting aside spin frames for just one minute}
\end{figure}

On the old palatial estate she was often found teasing highballs, with a gaggle of brittle men and women orbiting her, caught not in a net, but in her spell.

\end{minipage}
\newpage

\lettrine{N}ina was recalling earlier in the day, her agape at a purposeful murmuration of starlings, their synchronized flight of form and almost deliberate grace. That extraordinary improbability of such order emerging without individuated intention. She considered, what if these \emph{ert} birds were instead \emph{inert} leaves, aimlessly swept up by the wind and was minded in all her ertness of the Heegner numbers, those prime fragments aligning miraculously into order without apparent cause.

\begin{figure}[H]
\includegraphics[width=\linewidth]{Illustrations/vign-leafMurmer.png}
\caption{People pause to ponder the inert leaves expressing purpose}
\end{figure}

Nina was not drawn to the cosmological conjurings of her eminent academic father, Abhay,\footnote{A. Ashtekar, \emph{New Variables for Classical and Quantum Gravity}, Phys.~Rev.~Lett.~\textbf{57}, 2244–2247 (1986).}
 as all that space-time stuff left her a little cold. She was a hard nosed arithmetic number theorist staring at the number\footnote{all but twenty-six trillion, two hundred fifty-three billion, seven hundred forty-one million, two hundred sixty-four thousand, seventy-eight hundred twenty-five.} she had written on the back of a passing waiter who had laboured too long with the woman sat opposite, \(
e^{\pi \sqrt{163}}.
\)
Tantalizingly close to being an integer, yet fundamentally transcendent. \( \mathbb{Q}(\sqrt{-163}) \): the last and largest Heegner number for Ramanujan, the heir to a private joke the integers play with infinity.
 
\newpage

\textbf{Heegner Numbers:} special integers \( n \) for which the imaginary quadratic field \( \mathbb{Q}(\sqrt{-n}) \) has class number one (unique factorization in its ring of integers). There are \href{https://colab.research.google.com/drive/11T6wIoD9c4Wjrojnlqv2p7QSgzDllg44?usp=sharing}{seven Heegner numbers}
\[
n = 1, 2, 3, 7, 11, 19, 43, 67, 163.
\]

\noindent
As Nina reconjured the murmuration dissolving into randomness, she thought how the universe’s irrational swirls will, at times, resolve into something that looks whole: maths has a taste for such rare alignments. 

\begin{table}[ht]
\centering
\begin{tabular}{ccc}
\toprule
\( n \) & \( H_n \) & \( e^{\pi \sqrt{H_n}} \) \\
\midrule
1 & 1 & 23.140692632779 \\
2 & 2 & 85.019695223207 \\
3 & 3 & 230.7645883191458 \\
4 & 7 & 4071.932095225261 \\
5 & 11 & 33506.143065592438 \\
6 & 19 & 885479.777680154319 \\
7 & 43 & 884736743.999777466034906 \\
8 & 67 & 147197952743.999998662454224506 \\
9 & 163 & 262537412640768743.99999999999925007259719818568888 \\
\bottomrule
\end{tabular}
\caption{Tabulated list of \( n \), \( H_n \), and \( e^{\pi \sqrt{H_n}} \)}
\end{table}

\noindent
For \(H_7, H_8, H_9\) the values are “all but” integers, to increasingly absurd precision. And Nina pondered 23. Could it be mere coincidence that \(e^{H(1)\pi}\), with \(H(1)=\sqrt{1}\) the first Heegner number, presses in on \(\mathcal{L}=23\): the first integer whose root of unity fails to factor uniquely\footnote{See Arnie Heavyside for the full treatment of \emph{Larmor’s 23}, every self-respecting mathematician’s favourite number. All Heegner numbers are important in understanding complex multiplication of elliptic curves and class numbers of \emph{imaginary quadratic fields}.}, the very integer that sits like a keystone between two transcendental whispers?
\[
22.4592 < \pi^{e} \;<\; \mathcal{L} \;<\; e^{\pi} < 23.1407.
\]
The way \(23\) is bracketed: nudged from below by \(\pi^{e}\) and from above by \(e^{\pi}\) felt like a marking on the stone, a chiselled hint that some alignments are rare on purpose.


\medskip

Nina sat quietly, reflecting on the art of epigraphy\footnote {The deciphering of ancient inscriptions in bronze and granite.} and Euclid’s assertion: \textit{“The laws of nature are but the mathematical thoughts of God.”} Epigraphers may tease out stories etched in granite, but mathematicians no less read the faint striations in the primes: the hidden gyrations of order beneath the noise.


Her thoughts wandered to Stanislaw Ulam, a physicist better known for the Monte Carlo method and the mathematics of nonlinear systems. Ulam’s arrangement of numbers in spiraling epigraphic form, the runes of a lurking structure made manifest. 

\begin{figure}[H]
    \centering
    \includegraphics[width=\linewidth]{Illustrations/ulam411000.png}
    \caption{\href{https://colab.research.google.com/drive/1W_AsB1X2cMbZNaJTkTq3Vix2pAI6rMy9?usp=sharing}{Ulam Spiral} for Fermat and Gauss primes}
    \label{fig:Ulam411000}
\end{figure}

By plotting the natural numbers in a spiral and marking the primes in red, a striking affinity for their alignment along diagonals emerges. 

In the image below prime numbers are written in red $4n+1$ and blue $4n+3$ and their affinity to diagonals (almost as if the stone itself preferred straight cuts) is almost a cause for pause.

\begin{figure}[H]
    \centering
    \includegraphics[width=\linewidth]{Illustrations/Ulam-41.png}
    \caption{Ulam Spiral for Fermat and Gauss primes}
    \label{fig:Ulam41}
\end{figure}

The pattern appeared to Nina as almost deliberate, as if a spider counting in its native octal had spun them into place. And then it occurred to her. Why limit to a square lattice?
Nina then proceeded to \href{https://cyclotomic-ulam.vercel.app/}{code} up a Vercel-hosted app hoping it would reveal something distinct. Nina stared at the hexagonal field of numbers urging something to stir from its cold surface. The structuring was Eisenstein, a tessellation breathing through sixty degree articulations of symmetry rather than the orthogonality of the square. The two families of Fermat and Gauss held their own orbits. An arithmetic geography built from residues mod six, the home of primes shaped by the rhythm of the forms 6n plus 1 and 6n minus 1. The broad Fermat lineage, those primes of the form 4n plus 1, wandered between both mod six territories. The Gaussian primes, born of 4n plus 3, falling entirely into the 6n minus 1 class, a single red river threading the lattice. The pattern was not random. It whispered.
Then she changed the starting number to Gauss's 41:

\begin{figure}[h]
  \centering
  \includegraphics[width=\linewidth]{Illustrations/41EisensteinUlam.png}
  \caption{Ulam Spiral on an Eisenstein Lattice with blue Fermat primes.}
\end{figure}

The north and southwest rays shimmered with an illusion of purity. At first sight it appeared a river of Fermat ancestry, every prime dressed in the livery of four n plus one. The Gaussian lineage red and insistent goes East. She would write a \href{https://www.researchgate.net/publication/398122608_Numbers_That_Resolve_Primes}{paper} on this resolving of primes along the three principle axis of the hexagonal lattice.

\begin{tcolorbox}[title={Prime residues, principal rays, and $12 = \mathrm{lcm}(4,6)$}]
Primes larger than $3$ obey two simple congruence rules.

\begin{enumerate}
  \item Modulo $6$ every such prime lies in a $6n\pm1$ class:
  \[
    p>3 \quad\Rightarrow\quad p \equiv 1 \pmod{6} \quad\text{or}\quad p \equiv 5 \pmod{6}.
  \]

  \item Modulo $4$ every such prime lies in either the ``Fermat'' or the ``Gaussian'' family:
  \[
    p \equiv 1 \pmod{4} \qquad\text{(broad Fermat type, splits in $\mathbb{Z}[i]$)}
  \]
  \[
    p \equiv 3 \pmod{4} \qquad\text{(Gaussian type, inert in $\mathbb{Z}[i]$)}.
  \]
\end{enumerate}

To see how these families sit on an Eisenstein hexagonal lattice it is natural to combine both congruence conditions, so:
\[
  \mathrm{lcm}(4,6) = 12.
\]
Working modulo $12$ we obtain four allowed residue classes for primes larger than $3$:
\[
  p \equiv 1,5,7,11 \pmod{12}.
\]
Within these the two families separate cleanly:
\[
  \begin{array}{ccl}
    p \equiv 1 \pmod{4} &\Longleftrightarrow& p \equiv 1 \text{ or } 5 \pmod{12} \quad\text{(broad Fermat)}\\[2mm]
    p \equiv 3 \pmod{4} &\Longleftrightarrow& p \equiv 7 \text{ or } 11 \pmod{12} \quad\text{(Gaussian)}.
  \end{array}
\]
\end{tcolorbox}

\begin{tcolorbox}[title={Prime residues, principal rays, and $12 = \mathrm{lcm}(4,6)$}]
On the Eisenstein lattice we model points by Eisenstein integers
\[
  z = a + b\,\omega,\qquad \omega = -\tfrac12 + i\tfrac{\sqrt3}{2},\qquad \omega^{3}=1,\ \omega\neq1,
\]
which generate a hexagonal tiling with six principal rays.  
These rays are the lines
\[
  \{n\,\omega^{k} : n\in\mathbb{Z}_{\ge0}\},\qquad k=0,1,2,3,4,5,
\]
and in the usual integer labelling of the hex spiral each such ray corresponds to an arithmetic progression
\[
  p(n) = 3n + c_{k},
\]
for a suitable constant $c_{k}$ depending on the direction $k$.  
Hence primes on any principal ray arise from a single linear form in $n$ and so occupy at most two of the four allowed residue classes modulo $12$.  

In particular a given ray may appear visually dominated by broad Fermat primes ($4n+1$ types, residues $1$ or $5$ mod $12$) or by Gaussian primes ($4n+3$ types, residues $7$ or $11$ mod $12$).  
The apparent colour bias along a ray in the Eisenstein prime spiral is therefore a direct expression of this combined mod $4$ and mod $6$ structure compressed into the modulus $12$ and then projected along the linear forms $3n+c_{k}$ determined by the powers of $\omega$.
\end{tcolorbox}


\section*{Prime Discrimination}

The young waiter being a researcher in encryption key‑exchange protocols by day (a fact which spared him some castigation for her present neglect), was as good a muse as she would get today.

“Did you know,” Nina begins, stirring her highball and loosening her collar, “that Euler—or was it Gauss? ... toyed with a polynomial that spits out primes?
\[
p(n)=n^2+n+41.
\]
Primes for \(n=0\) through \(40\); then at \(n=41\), the spell breaks.”

\begin{figure}[H]
\includegraphics[width=\linewidth]{Illustrations/vign-ninaCafe.png}
\caption{Nina setting aside spin frames for just one minute}
\end{figure}

The waiter posturing appropriately, “A polynomial generating primes? That sounds almost algorithmic for something said to be random.”

“It is—fleetingly,” Nina said, eyes bright. “Look at the discriminant.”
\[
\Delta=b^2-4ac=(1)^2-4(1)(41)=-163.
\]
“Not just any negative: the very \(-163\) that rules the exquisite modularities of \( \mathbb{Q}(\sqrt{-163}) \). It’s no accident that \(e^{\pi\sqrt{163}}\) brushes against an integer; that’s the symmetry talking.”



The young waiter who was a researcher in all matters encryption, digital signatures, authentication and key exchange protocols and the likes - a fact that had saved him from some castigation for her present neglect- was as good a muse as she was to get today to put her idlings to.

"Did you know," Nina begins, as she stirs her high ball undoing her collar, "that Euler—or was it Gauss?—once played with a polynomial that generated primes? \(n^2 + n + 41\). It produced primes for \(n = 0\) to \(40\), but then at \(n = 41\), the magic stopped."

The young waiter, interest piqued "A polynomial generating primes? That can be right almost... algorithmic when they are random?"

Nina nodding, her eyes now a sparkling. "It is, in a way. But it’s also fleeting. At \(n = 41\), the output becomes composite, and the spell is broken. It’s like quadratic reciprocity in which the hidden symmetries of numbers suddenly show themselves, as if reminding us they’re in control."

The waiter grinning, "Quadratic reciprocity ... now that’s a backbone for encryption. But what about primes like 163? I’ve heard it’s... special."

Just at that moment, the rather elegant lady that the waiter had been over attentive with, pipes up "All that talk of quadratics; all that school stuff gives me the willies. I thought you two were sharing trade secrets on the back of the napkin the way you were carrying on."

Nina, not picking up on the her tone, adopts the patient tone of someone explaining a cherished concept. The young waiter listens politely.

"For your one, let’s revisit Euler’s polynomial," Nina began, "the one I mentioned earlier, \(p(n) = n^2 + n + 41\). This polynomial generates primes for \(n = 0\) through \(40\), but at \(n = 41\), it fails spectacularly, producing a composite. Why does this happen?" She reached for her notebook and wrote:
\[
p(n) = n^2 + n + 41.
\]
"Now," she continues, "to understand its properties, we examine your discriminant ... what's your name, "Alma". OK Alma. You'll surely remember the coefficients of the general quadratic equation \(ax^2 + bx + c = 0\). Fpr our Euler they are, \(a = 1\), \(b = 1\), and \(c = 41\). The <<discriminant, \(\Delta\), is calculated as:
\[
\Delta = b^2 - 4ac.
\]"
With a flourish, she writes:
\[
\Delta = (1)^2 - 4 \cdot 1 \cdot 41 = 1 - 164 = -163.
\]

"Do you see that?" she asked, peering over her glasses. "The discriminant, \(\Delta = -163\), is not just a random number ... it’s deeply tied to the polynomial’s prime-generating behavior. It connects to the modular properties of quadratic fields, specifically those involving \(\sqrt{-163}\)."

% \subsection*{A Common Discriminant: Helgott’s Formula}

Nina gesturing for the waiter to get another High ball, "Now, let me introduce you, Alma, to another remarkable quadratic, this one from Helgott: \(q(n) = n^2 - 79n + 1601\). Like Euler’s formula, it generates primes for consecutive values of \(n\). Let’s calculate its discriminant." She scribbled:
\[
q(n) = n^2 - 79n + 1601.
\]
"Here, \(a = 1\), \(b = -79\), and \(c = 1601\). Applying the same formula:
\[
\Delta = (-79)^2 - 4 \cdot 1 \cdot 1601 = 6241 - 6404 = -163.
\]"
Alma’s eyebrows archinh, "Wait, both polynomials have the same discriminant?" Nina Sketched the quadratics to nail home the point.

\begin{figure}[H]\centering
\includegraphics[width=0.8\textwidth]{Illustrations/HelgotEuler.png}
\caption{\href{https://colab.research.google.com/drive/1hxhBZ2Ns7NXTzxHeYBMa-W1F4MXvtmEB?usp=sharing}{Quadratic Generators} of Primes.}
\label{HelgotEuler:fig}
\end{figure}

"Exactly," Nina  self-satisfied smiling having breached the defences of a truculent student. "Both Euler’s and Helgott’s quadratic formulas share the discriminant \(\Delta = -163\). But the connection runs deeper. This discriminant links to the exceptional modular properties of the number field \(Q(\sqrt{-163})\). It’s no coincidence that \(e^{\pi \sqrt{163}}\) is almost an integer as it’s a reflection of these underlying symmetries." Alma looking phased and dazed was calling on another waiter to make up her gin quota for the day. Plainly peaked and Nina needed a more attentive audience. The waiter dutifully arrives.

"Now, here’s where it gets fascinating," Nina continuing as if he had never left. "No other quadratic formula with a discriminant of \(-163\) generates prime numbers for consecutive positive integer values of \(n\). The rarity of \(-163\) makes it unique; as if nature is choosing one number to weave together primality, modularity, and geometry."

Our dutiful waiter obliges, "So, is this where the connection to primes ends?"

"Not even close," Nina said. "Ramanujan might have called it a 'bridge to transcendence.' It’s tied to the \href{https://colab.research.google.com/drive/1K3itZdpqQkzb5UDWO-3MrIe6-KwopK0q?usp=sharing}{near-integer} property of \(e^{\pi \sqrt{163}}\):
\[
\mathcal{R} \approx 262537412640768743.99999999999925.
\]
Two transcendentals—\(e\) and \(\pi\)—and a quadratic irrational \(\sqrt{163}\) combine to create something almost perfect."

Almost conspiratorial, but mostly incoherently so "And here’s a detail almost nobody talks about—this number is absurdly close to:
\[
\mathcal{R}_p\approx M = 640320^3 + 744= 262537412640768744.
\]
Through iterative \href{https://colab.research.google.com/drive/1hn_NBpzrlv1TNKeRQAPDLabW04Jya0tv?usp=sharing}{Shor}-style analysis and \href{https://colab.research.google.com/drive/1Io6trNJBcgsMHd2it5qg-_ZYuJLJOb8p?usp=sharing}{classical factorization}, we find:
\[
M = 2^3 \cdot 3 \cdot 10939058860032031,
\]
for a prime, \(10939058860032031\). Raising \(\mathcal{R}\) to nth powers and we have what Demi Moorish might refer to as an nth power puboid but for the cubic two:
\begin{align*}
M^n &= (2^3 \cdot 3 \cdot 10939058860032031)^n\\
&= 2^{3n} \cdot 3^n \cdot 10939058860032031^n
\end{align*}
These powers of M deliver the integers in the following table.

\begin{longtable}{|c|p{12cm}|}
\hline
\textbf{n} & \textbf{Integer Part of \( \mathcal{R}^n \)} \\
\hline
1 & 262537412640768743 \\
2 & 68925893036109279891085639286943729 \\
3 & 18095625621654510801615355531263439695149914167019841 \\
4 & 4750778730825177725463920948909721363447613520911266766459844650576849 \\
\hline
\end{longtable}
\noindent
But considering \(\mathcal{R}\) in full we see that the near-integer \href{https://ramanunjan-rhapsody.vercel.app/}{magic} slips away as \(n\) grows—the decimal tails start to wander:

\begin{longtable}{|c|p{12cm}|}
\hline
\textbf{n} & \textbf{Decimal Part of \( \mathcal{R}^n \)} \\
\hline
1 & 0.62537412640768743999999999999250072597198185688879353856337336990862 \\
2 & 0.8925893036109279891085639286943768000000000163738644209234607572328 \\
3 & 0.0956256216545108016153555312634547066300647710749759999999901236935249 \\
4 & 0.5077873082517772546392094890972661821449171803947136631874740636874078\\
\hline
\end{longtable}

The runs of 9s and 0s tell the story of that fading precision:
\vspace{0.3cm}

\begin{tabular}{|c|c|c|}
\hline
\textbf{n} & \textbf{Longest Run of 9s} & \textbf{Longest Run of 0s} \\
\hline
1 & 12 & 2 \\
2 & 1 & 9 \\
3 & 8 & 2 \\
4 & 1 & 1 \\
\hline
\end{tabular}
\vspace{0.3cm}

"See? At \(n=1\), it’s freakishly precise such that by \(n=4\), the perfection’s gone as the fog lifts from your hand."

Nodding knowingly, the waiter knowingly says: "It sounds modular—elliptic curves, symmetry?"

"Exactly!" Nina openly grinning. "It’s a singularity: a moment where infinite expansions converge into an almost-perfect whole."

% \subsection*{The Continued Fraction}

Nina, after taking on an iced-tea flattener, begins to calculate the \href{https://colab.research.google.com/drive/1gHUha6vLc_Mg0x224fmExqYGg0-iLQwP?usp=sharing}{continued fraction} representation of $\mathcal{R}$:
\begin{align*}
\mathcal{R} &\approx 262537412640768743.9999999999992500725971981856\ldots\\
&= [\,262537412640768743; 1, 1333462407511, 1, 8, 1, 1, 5, 1, 4, \dots\,].
\end{align*}
Talking to herself now that the waiter has been called away "The enormous second partial quotient appears because $\mathcal{R}$ is only $7.499\times 10^{-13}$ less than an integer, so that
\[
\mathcal{R} = N + (1 - \varepsilon), \quad \varepsilon \approx 7.499\times 10^{-13},
\]
and thus the reciprocal $1/\varepsilon$ produces the large term $1333462407511$. In stark contrast, the golden ratio $\varphi = [1;1,1,1,\dots]$ has all partial quotients equal to $1$, giving it the slowest possible rational approximations and earning it the title of the <most irrational> number. By the same playful logic, $\mathcal{R}$, with its colossal early term, sits at the opposite extreme ... among the <most rational irrationals.>"

Alma, having overheard, drifted over. Not herself that well versed in matters of the attentive mind despite a life-time of being partially self-aware she was strangely drawn into their other worldly fascinations.

Nina looking up from her gin, thinks better of it, and flips a page in her magazine. Nina then thought of her Uncle Abhay’s love of physics. 

"To detect a ripple in spacetime, the precision required was on the order of \(10^{-21}\)—an astonishing accuracy. About fifteen terms in the continued fraction would do it. Beyond that, the error becomes negligible."

Looking it up on the LIGO site, she reads that the detectors are sensitive to gravitational-wave \emph{strain} at the level of about $10^{-21}$,\footnote{\href{https://dcc.ligo.org/LIGO-P070082/public}{LIGO overview: “sensitive to gravitational wave strains smaller than one part in $10^{21}$”}} 
with path-length changes even smaller than a proton’s width.\footnote{\href{https://www.ligo.caltech.edu/page/detection}{LIGO public page: detecting changes $\sim \tfrac{1}{1000}$–$\tfrac{1}{10000}$ of a proton’s width}}

“Such feats lean on constants like $\pi$,” she thought. “A sizable error in $\pi$ would wreck that precision.”

Yet for
\[
\mathcal{R}=e^{\pi\sqrt{163}}=N-(\varepsilon),\qquad 
\varepsilon \approx 7.499274028\times 10^{-13},
\]
the \emph{relative} miss is
\[
\frac{\varepsilon}{\mathcal{R}}
\;=\;
\frac{7.499274028\times10^{-13}}{2.62537412640768744\times10^{17}}
\;\approx\;
2.856459181\times 10^{-30}.
\]
Compared to LIGO’s characteristic strain scale $h\sim 10^{-21}$, this is tinier by a factor of
\[
\frac{10^{-21}}{\varepsilon/\mathcal{R}}
\;\approx\;
3.50\times 10^{8},
\]
% i.e.\ about \(\mathbf{3.5\times 10^8}\) times smaller in fractional terms. 

“Even our universe-to-proton analogies feel blunt,” Nina smirks. “\(\mathcal{R}\)’s nearness to an integer beats LIGO’s fractional sensitivity by nearly nine orders of magnitude.” 

"Touche-eh," Alma replied.

Nina closes her notebook looking at the equation once more. 

\subsection*{Narkilar and Other Dangerous Sleeps}

Nina blamed her father for this sort of thinking. Abhay had a diddle-squat interest in the numerological borderlands of physics: the ones where Narkilar’s\footnote{Jayant V. Narlikar, in his classic text \emph{An Introduction to Cosmology}, often explored numerological coincidences and large dimensionless ratios in physics \cite{narlikar2002}.}.
 tea-stained conjectures about “natural scales” rubbed elbows with Dirac’s Large Number Hypothesis\footnote{P.A.M. Dirac, ``A new basis for cosmology,'' \emph{Proceedings of the Royal Society of London. Series A}, vol. 165, no. 921, pp. 199–208, 1938.}. A graveyard of otherwise respectable intellects, he would joke, but still a place he enjoyed strolling through. 

\medskip

\noindent So perhaps it was inevitable that she’d idly notice, between sips, that the reciprocal of the squared Heegner value,
\[
\frac{1}{\mathcal{R}^2} = \frac{1}{6.89258\times 10^{34}} \;=\; 1.4508\times 10^{-35},
\]
lands squarely in the Planckish neighbourhood — \(\hbar = h/2\pi \approx 1.0545\times 10^{-34}\ \mathrm{Js}\).

Abu, Abhay would have said, “Of course it does as all roads in the zoo of dimensional analysis lead back to the same cages. But look there at those Joule Seconds ... so \href{https://www.researchgate.net/publication/390318980_Redefining_Planck's_Constant_via_Cosmological_Time_Ramanujan_Time_and_the_Wien_Displacement_Field}{anthropic}.” 

% >>>>>> ANNEX C (cut-sheet v2): the syntaxology of mass scales and Planck--Einstein symbolic regression — block commented out below, pointer left in text <<<<<<
\textit{(Nina's game of deriving Planck and Einstein by symbolic regression is played out in full in Annexe C.)}

% \subsection*{Syntaxology of Natural Mass Scales}
% 
% Physics often advances by staging a duel between competing metaphors: waves vs.\ particles, geometry vs.\ inertia, localisation vs.\ collapse.  
% Sometimes these duels produce balances that are “pre-Einstein compatible”: simple, physically motivated constructions that already echo later theory.  
% 
% Here is one such duel, in which the Quantum that is Nina says: “this is how small you can get,” and gravity replies: “this is how big you must be.” Behold the Compton wavelength,
% \[
% \lambda_C=\frac{h}{mc},
% \]
% capturing how tightly a particle of mass \(m\) can be pinned down before quantum fuzziness sets in while the Schwarzschild radius,
% \[
% r_s=\frac{2Gm}{c^2},
% \]
% measures how much mass bends spacetime before it becomes a black hole. 
% Now if we must insist that they speak to each other,
% \(
% \lambda_C=r_s,
% \)
% the solution is the \emph{Planck mass}:
% \[
% m_P=\sqrt{\frac{hc}{G}} \;\;\sim 5.456\times10^{-8}\,\mathrm{kg},
% \]
% marking the point where quantum localisation and gravitational collapse \emph{agree on the rules}.  
% In the same spirit, as we will see in a minute, a Newtonian universe mass,
% \[
% M_U\sim \frac{c^3}{GH},
% \]
% can fall out from nothing more than binding energy plus a cosmic clock.
% 
% Nina's dad had once shown her a longer list of “natural masses” \href{https://mcmurmerings.wordpress.com/2021/08/03/a-complete-set-of-fundamental-mass-combinations/}{cooked up} from the constants \(c, h, G, H\) by nothing more exotic than unit juggling.
% 
% \href{https://weinberg-mass-hierarchy.vercel.app/}{Dimensional analysis} asks us to combine these constants so that the units cancel into kilograms,  
% \[
% m \sim c^\alpha h^\beta H^\gamma G^\delta,
% \]
% with the exponents fixed by three unit equations \([M],[L],[T]\).  
% There are four independent solutions:  
% \[
% \begin{aligned}
% m_1 &= \frac{c^3}{GH} \;\;\sim 1.8\times10^{53}\ \mathrm{kg}, & &\text{mass of the Hubble sphere},\\[6pt]
% m_2 &= \frac{hH}{c^2} \;\;\sim 1.6\times10^{-68}\ \mathrm{kg}, & &\text{a graviton/self–gravity scale},\\[6pt]
% m_3 &= \Big(\tfrac{h^3 H}{G^2}\Big)^{1/5} \;\;\sim 4.3\times10^{-20}\ \mathrm{kg}, & &\text{an intermediate curiosity},\\[6pt]
% m_4 &= \Big(\tfrac{hc}{G}\Big)^{1/2} \;\;\sim 5.456\times10^{-8}\ \mathrm{kg}, & &\text{the Planck mass}.
% \end{aligned}
% \]
% 
% These seeds already span more than \(10^{121}\) in magnitude: the same vast gulf between quantum field theory’s vacuum energy and the observed dark-energy density.  
% It was this kind of gap Eleanor and Jovannah hoped to bridge by trading the infinities of a continuous manifold for the cleaner bookkeeping of a lattice, so sidestepping the arbitrary ansätze of \href{https://www.researchgate.net/publication/23671110_Regularization_Renormalization_and_Dimensional_Analysis_DimensionalRegularization_meets_Freshman_EM}{\emph{regularisation} and \emph{renormalisation}}
% 
% , those conjuring acts where, if you allow water to turn into wine, anything else, even the charge on an electron, is possible on the menu.  
% 
% Once we have the seeds \(\{m_1,m_2,m_3,m_4\}\), we can generate further scales by taking geometric means:
% \[
% m_{ij}=\sqrt{m_i m_j}, \qquad m_{ijk}=(m_i m_j m_k)^{1/3}, \ldots
% \]
% This combinatorics produces some intriguing echoes:  
% the Planck mass itself reappears as \(m_{12}\);  
% \(m_{24}\sim 3\times10^{-38}\,\mathrm{kg}\), suggestive of neutrino bounds;  
% \(m_{234}\sim 3.4\times10^{-32}\,\mathrm{kg}\), in neighbourhood of electron mass.  
% 
% It was this proliferation that inspired Weinberg’s image of a “dial”: spin it one way, the crossings fall near mesons; turn it another, and they brush past neutrinos or electrons.  
% The pattern is seductive: real particles seem to lurk near these geometric means.  
% It looks elegant. But elegance is not explanation. Unless some principle fixes the dial at a definite setting, most of these constructions remain suggestive curiosities rather than physics. When formulas from different corners of physics suddenly rhyme we might call that \emph{syntaxology}.
% 
% \subsection*{Symbolic Regression Without a Map}
% In the spirit of an unguided machine we can \href{https://colab.research.google.com/drive/1Jc4zmPlcXqZImYC3wqJJT7Ry__WeUmSY?usp=sharing}{treat} each mass expression as a \emph{syntactic object} and quantify its \emph{structural complexity}.  
% A simple industry-standard proxy is the \emph{functional form depth} which comprises the maximum nesting level of operators and functions in the expression tree.\footnote{This is essentially the tree-depth measure used in symbolic regression fitness functions.  
% For example, $\sqrt{a\,b}$ has depth $2$ (multiply, then root), whereas $\log\!\sqrt{a\,b}$ has depth $3$.  
% This is a coarse but widely-used measure of algebraic complexity.}
% 
% \begin{center}
% \begin{tabular}{lcc}
% \hline
% Expression & log$_{10}$ mass (kg) & Complexity depth \\
% \hline
% $hH/c^2$ & $-68.4$ & 2 \\
% $h/(c\,R_U)$ & $-67.6$ & 3 \\
% $\sqrt{hH/G}$ & $-59.9$ & 2 \\
% $h/(m_p c)$ & $-27.3$ & 3 \\
% $m_P=\sqrt{h c/G}$ & $-7.6$ & 2 \\
% $c^3/(GH)$ & $52.0$ & 2 \\
% $\sqrt{m_P M_U}$ & $22.2$ & 3 \\
% \hline
% \end{tabular}
% \end{center}
% 
% Here the $y$-axis is $\log_{10}(\mathrm{mass/kg})$, and the $x$-axis is \emph{syntactic complexity depth}.  
% Even without a physical story, the low-complexity forms cluster at physically meaningful scales (Planck, baryonic universe, proton scale), while arbitrary composites wander.
% 
% \begin{figure}[h]
% \centering
% \includegraphics[width=0.9\textwidth]{Illustrations/mass_vs_complexity.png}
% \caption{Mass scales (log$_{10}$ kg) vs.\ syntactic complexity depth.  
% Lower complexity often aligns with physically motivated scales; higher complexity contrived numerology.}
% \end{figure}
% 
% This “complexity spectrum” offers a crude visual heuristic:  
% \emph{physics lives where syntax is simple and the constants are there for a reason}. The space of \emph{allowable} forms is infinite; physical reality occupies a vanishingly small subset.
% 
% 
% 
% To move from rhyme to reason, one must press these coincidences into \href{https://symbolic-regression.vercel.app/}{models}. Only then can we see whether they carry explanatory weight once embedded in dynamics. 
% 
% The best way to see the difference is to build the simplest garden-variety model and watch where it succeeds and where it slips. 
% 
% Let that garden be Newton’s.  Picture the universe as he sees it, as nothing more than a sphere of dust: matter with no pressure, expanding or contracting with time. 
% Give it a radius \(R(t)\), a density \(\rho(t)\), and an expansion rate given by a "Hubble parameter":
% \[
% H(t)=\frac{\dot R}{R}.
% \]
% 
% Now imagine standing on the boundary and placing a test particle there.  
% Newton tells us to balance its kinetic energy against the pull of the whole sphere inside. 
% That energy statement is
% \[
% \frac{1}{2}\dot R^2 - \frac{GM}{R} = \frac{k c^2}{2},
% \qquad
% M=\frac{4\pi}{3}\rho R^3,
% \]
% where the extra term \(k c^2/2\) encodes whether the geometry is open, flat, or closed.  With a little algebra\footnote{By multiply both sides by \(\displaystyle \frac{2}{R^2}\) and substituting for \(M\):
% \[
% \frac{\dot R^2}{R^2}-\frac{2G}{R^3}\!\left(\frac{4\pi}{3}\rho R^3\right)+\frac{k c^2}{R^2}=0
% \;\;\Longrightarrow\;\;
% H^2-\frac{8\pi G}{3}\rho+\frac{k c^2}{R^2}=0,
% \]}, the Newtonian energy balance reshuffles into
% \[
% H^2=\frac{8\pi G}{3}\rho-\frac{k c^2}{R^2}.
% \]
% The left-hand side is simply the square of the expansion rate. 
% On the right, we have a term proportional to the density and a term that looks like a curvature correction, with \(k\) playing the role of spatial geometry: positive for closed, zero for flat, negative for open. In other words, by juggling Newton’s kinetic and potential energies for a dust sphere, we stumble onto a formula that already carries the skeleton of Einstein’s cosmology. 
% 
% From our Newtonian form choose the spatially \emph{flat} case \(k=0\). In this limit the expansion is
% balanced by a specific density, the \emph{critical density},
% \[
% \rho_c \;\equiv\; \frac{3H^2}{8\pi G}.
% \]
% \noindent
% If we now take a Hubble-sized patch with radius \(R\sim c/H\), its mass is
% \[
% M_U \;\sim\; \rho_c\,\frac{4\pi}{3}R^3
% \;\sim\; \Big(\frac{3H^2}{8\pi G}\Big)\,\frac{4\pi}{3}\,\Big(\frac{c}{H}\Big)^3
% \;\sim\; \frac{c^3}{2G H}.
% \]
% This gives us a compact estimate for the mass within a Hubble volume. 
% It has the right shape, even the right order of magnitude. 
% But notice a hitch. If instead we try the different Newtonian shortcut of equating the Hubble recession speed $HR$ to the Newtonian escape velocity $\sqrt{2GM/R}$, we find:
% \[
% \rho \sim \frac{3H^2}{4\pi G},
% \]
% which is a factor of two larger than the correct critical density.  So even in Newton’s tidy garden, syntax needs a steady hand. The forms look familiar, but without Einstein’s pruning, the details slip.
%  
% 
% \newpage
% ===========================
% 
% \section*{Planck \& Einstein as directed symbolic regressors}
% 
% \subsection*{Planck’s Scaffolding}
% 
% At the turn of the twentieth century, Planck faced a hopeless counting problem: how many ways can energy be distributed among oscillators if energy is continuous? The answer was infinitely many—too many for statistical mechanics to function. So, he did what mathematicians often do: he added scaffolding.
% 
% He discretized. By chopping energy into chunks of size \(\epsilon\), the infinite became finite. With allowed energy levels \(0, \epsilon, 2\epsilon, \dots\), the combinatorics became tractable, entropy followed from Boltzmann’s formula, and thermodynamic laws yielded a plausible energy distribution. The plan was to solve the finite case, then let \(\epsilon \to 0\) to recover the classical world.
% 
% But when Planck removed the scaffolding and took the limit, the formula collapsed into the Rayleigh–Jeans catastrophe. Only by keeping \(\epsilon = h\nu\) finite did the theory agree with observation. What began as an artificial regularization had become the law itself.
% 
% This is symbolic regression at its most treacherous: introduce a fiction to make the math work, then discover that removing it breaks reality. Planck’s “temporary” discretization turned out to be nature’s permanent constraint. The scaffolding was the physics.
% 
% Mathematics offers an infinitude of dimensionally consistent forms.  
% Matter prunes these to what the universe actually contains.  
% In Penrose’s tripartite framing, it takes a Mind of the form and complexity of Einstein's to supply the guiding hand of heuristics and constraints. Seen through a Machine Learning perspective, Einstein's uncovering of general relativity was by \emph{guided symbolic regression}:
% \[
% \text{terminals: } (g_{\mu\nu},\,T_{\mu\nu},\,G,\,c),\quad
% \text{operators: } (\nabla,\,\text{contraction},\,\text{variation}),
% \]
% with strict \emph{fitness constraints}:
% general covariance, equivalence principle, Newtonian limit, energy–momentum conservation. From countless syntactic possibilities, only his mighty equation survives,
% \[
% G_{\mu\nu}=\frac{8\pi G}{c^4}T_{\mu\nu}.
% \]
% 
% \subsection*{Principle and Pruning in the Landscape of Syntax}
% 
% Planck and Einstein were, in essence, directed symbolic regressors: explorers of syntax guided by principle. Where mathematics offers an infinite landscape of sterile coincidences, they pruned with constraints sharp enough to reveal necessity. Physics resides only in those rare intersections where form meets principle. From Newton’s garden to Planck’s scaffolding to Einstein’s symmetries, the lesson repeats: without constraint, algebra blooms into illusion; with it, syntax yields law. The constants of nature are not curiosities but tokens in the grammar of reality, anchors in a language where simplicity signals truth, and where only a mind, human or machine, can keep invention from drifting away from discovery.
% 
%   
% Whether in 1905 or in the era of machine learning, the danger is the same: without principled constraints, your symbolic regression\footnote{Symbolic regression: a branch of machine learning which searches over a vast space of algebraic expressions to fit data. Without strong priors such as dimensional analysis, invariance principles, or known symmetries, it will happily produce elegant but meaningless formulae.} will return a forest of plausible nonsense.  
% 
% \begin{quote}
% On syntax, regression, and the grammar of reality:  
% 
% the regressive eye sees countless forms,  
% 
% the progressive mind prunes them to one.  
% 
% Syntax may wander, but truth lives  
% 
% where simplicity and principle meet.
% \end{quote}
% 
% Nina knew her whole scheme was half stagecraft, half physics, and maybe that was the charm.  But once you’d seen one “coincidence” like \(\mathcal{R}^2\) sidling up to \(\hbar\), you started spotting them everywhere, 
% like an earworm you couldn’t shake. And just then the waiter returned with another highball. A nod was as good as a wink. Sipping her third, eyes sparkling, she did not let him go. He is smiling (or is it smirking?) as he places the drink gently now beside her.
% 
% \begin{itemize}[leftmargin=1.5cm]
% 
% \item[Nina:] \textit{} You know, we’ve got it all wrong: anchoring time to caesium atoms. 
% What if we ditched that whole atomic-clock business?
% 
% \item[Waiter:] \textit{} And replace it with what, Professor?
% 
% \item[Nina:] The Cosmic Microwave Background. Its Wien peak today is about \(160.3~\mathrm{GHz}\). 
% Now measure that against the \emph{Planck time},
% \[
% t_P=\sqrt{\frac{\hbar G}{c^5}} \;\approx\; 5.39\times10^{-44}\,\mathrm{s}.
% \]
% Multiply the CMB peak frequency by this tick, and you get a pure, dimensionless number:
% \[
% \chi_0=\nu_{ CMB,0}\,t_P \;\approx\; 8.6\times 10^{-33}.
% \]
% 
% \item[Waiter:] A dimensionless “cosmic clock ratio,” then. But does it actually \emph{say} anything?
% 
% \item[Nina:] That’s the trick. Two ladders appear:
% 
% \emph{First, the static ladder (today):}
% \[
% \frac{L_H}{\ell_P} \;\approx\;\frac{t_H}{t_P} \;\approx\;\mathcal{R}^{7/2}.
% \]
% Right now, the Hubble scale measured in Planck units sits almost exactly on the half-integer rung \(7/2\).
% 
% \emph{Second, the dynamic ladder (future):}
% as the universe expands, \(\nu_{CMB}(t)\) redshifts and so does \(\nu_{CMB}(t)\,t_P\). 
% In about \(10^{11}\) years—when the CMB cools to \(\sim 4.6\) mK—this ratio will slide down to
% \[
% \nu_{CMB}(t)\,t_P = \mathcal{R}^{-2}.
% \]
% 
% \item[Waiter:] So the Planck time sets the ruler, the CMB sets the tick, and Ramanujan supplies the rungs?
% 
% \item[Nina:] Exactly. One coincidence frozen today at \(\mathcal{R}^{7/2}\), another awaiting us in the far future at \(\mathcal{R}^{-2}\). 
% Not numerology, but syntax.
% 
% \item[Waiter:] \textit{(smiles)} Syntax I can live with. Numerology I only drink with. 
% Just don’t try to pay the bill in Ramanujans. The register only takes rupees.
% 
% \end{itemize}
% 
% \footnotetext{%
% \textbf{Symbolic regression as Machine Learning):} search over equation structures (sums, products, powers, logs, compositions) to fit data; 
% without constraints (dimensions, priors, symmetries), it will happily overfit with mathematically impeccable yet physically pointless forms.}
% 
%%%%%%%%%%%%%%%%%%%%%%
% \section*{Mathematical Postscript: ramifications}
% 
% The first taste of prime–splitting comes not in any exotic number ring, but in the integers themselves.  
% Fermat’s theorem on sums of two squares says that an odd prime $p$ can be written as $p=a^2+b^2$  
% iff $p\equiv 1 \pmod{4}$. Already here, the rational primes fall into two infinite families:  
% those with $p\equiv 1 \pmod{4}$ that appear as sums of two squares, and those with $p\equiv 3 \pmod{4}$ that do not. Viewed through $\mathbb{Z}[i]$ (the Gaussian integers), this is simply the splitting law:
% \begin{align*}
% p \ \text{splits in }\mathbb{Z}[i] & \iff p\equiv 1 \pmod{4}, \\
% p \ \text{inert} & \iff p\equiv 3 \pmod{4}, 
% \quad 2 \ \text{ramifies}.
% \end{align*}
% 
% From here the pattern generalises. For any quadratic field $K=\mathbb{Q}(\sqrt{m})$, $m$ square-free, whether a prime $p$ \emph{splits}
% is determined by whether the field’s discriminant\footnote{For $ax^2+bx+c=0$ the familiar $b^2-4ac$ is a \emph{discriminant}. 
% In quadratic fields, the \emph{field discriminant} $\Delta_K$ plays an analogous role for arithmetic: it encodes the structure of 
% $\mathcal{O}_K$ and governs how rational primes factor—splitting, inertness, or ramification—inside $\mathcal{O}_K$.} $\Delta_K$ is a quadratic residue modulo $p$,
% and quadratic reciprocity turns this into explicit congruence conditions. That is why $-163$ is remarkable: it is the last imaginary
% quadratic field with class number $1$, so unique factorisation holds. By contrast, $-7$ is also a Heegner number, hence class number $1$
% as well, but primes still split there—class number $1$ guarantees unique factorisation, not the absence of splitting.
% 
% \paragraph{Orders and discriminants.}
% For a quadratic field $K = \mathbb{Q}(\sqrt{m})$, $m$ squarefree, the maximal order $\mathcal{O}_K$ and discriminant $\Delta_K$ are
% \[
% \mathcal{O}_K =
% \begin{cases}
% \mathbb{Z}\!\left[\dfrac{1 + \sqrt{m}}{2}\right], & m \equiv 1 \pmod{4},\ \ \Delta_K = m,\\[0.8em]
% \mathbb{Z}[\sqrt{m}], & m \equiv 2,3 \pmod{4},\ \ \Delta_K = 4m.
% \end{cases}
% \]
% 
% \paragraph{Two contrasting examples.}
% \[
% -5 \equiv 3 \pmod{4}
% \quad\Rightarrow\quad 
% \mathcal{O}_{\mathbb{Q}(\sqrt{-5})} = \mathbb{Z}[\sqrt{-5}], \quad \Delta_K = -20,
% \]
% \[
% -7 \equiv 1 \pmod{4}
% \quad\Rightarrow\quad 
% \mathcal{O}_{\mathbb{Q}(\sqrt{-7})} = \mathbb{Z}\!\left[\frac{1+\sqrt{-7}}{2}\right], \quad \Delta_K = -7.
% \]
% 
% \begin{table}[h!]
% \centering
% \renewcommand{\arraystretch}{1.25}
% \begin{tabular}{|c|c|c|c|p{5.0cm}|}
% \hline
% $m$ & $m \bmod 4$ & $\mathcal{O}_K$ & $\Delta_K$ & Splitting condition for $p \nmid \Delta_K$ \\
% \hline
% $-5$ & $3$ &
% $\mathbb{Z}[\sqrt{-5}]$ &
% $-20$ &
% $p$ splits $\iff \left(\frac{-5}{p}\right)=+1$,
% i.e.\ $p\equiv 1,3,7,9 \pmod{20}$;  
% inert otherwise; $2$ ramifies. \\
% \hline
% $-7$ & $1$ &
% $\mathbb{Z}\!\left[\frac{1+\sqrt{-7}}{2}\right]$ &
% $-7$ &
% $p$ splits $\iff \left(\frac{-7}{p}\right)=+1$,
% i.e.\ $p\equiv 1,2,4 \pmod{7}$;  
% inert otherwise; $7$ ramifies. \\
% \hline
% \end{tabular}
% \caption{How $m \bmod 4$ fixes $\mathcal{O}_K$, $\Delta_K$, and the prime-splitting law.}
% \end{table}
% 
% \paragraph{Small-prime illustration.}
% For $\mathbb{Q}(\sqrt{-5})$:
% \[
% \text{Split: } 3,7,13,17,23, \dots \qquad \text{Inert: } 11,19,29, \dots
% \]
% A concrete split in elements:
% \[
% (3)=(1+\sqrt{-5})(1-\sqrt{-5}).
% \]
% 
% For $\mathbb{Q}(\sqrt{-7})$ (write $\omega=\tfrac{1+\sqrt{-7}}{2}$ so $N(a+b\omega)=a^2+ab+2b^2$):
% \[
% \text{Split: } 2,11,23,29,37, \dots \qquad \text{Inert: } 3,5,13,17,19, \dots
% \]
% A concrete split:
% \[
% (2)=(\omega)(\omega'),\qquad \omega'=\tfrac{1-\sqrt{-7}}{2},\qquad N(\omega)=2.
% \]
% 
% \paragraph{Heegner fields still split.}
% All nine imaginary quadratic fields with class number $1$ (the Heegner list) are UFDs, yet each has infinitely many split primes.  
% For $K=\mathbb{Q}(\sqrt{m})$ with Heegner $m$, the split primes are those with $\left(\frac{\Delta_K}{p}\right)=+1$:
% 
% \begin{table}[h!]
% \centering
% \renewcommand{\arraystretch}{1.15}
% \begin{tabular}{|c|c|p{10cm}|}
% \hline
% $m$ & $\Delta_K$  & Split congruences for $p\nmid \Delta_K$ \\
% \hline
% $-1$  & $-4$   & $p \equiv 1 \pmod{4}$. \\
% $-2$  & $-8$   & $p \equiv 1,3 \pmod{8}$. \\
% $-3$  & $-3$   & $p \equiv 1 \pmod{3}$. \\
% $-7$  & $-7$   & $p \equiv 1,2,4 \pmod{7}$. \\
% $-11$ & $-11$  & $p \equiv 1,3,4,5,9 \pmod{11}$. \\
% $-19$ & $-19$  & $p \equiv 1,4,5,6,7,9,11,16,17 \pmod{19}$. \\
% $-43$ & $-43$  & $\left(\frac{-43}{p}\right)=+1$ (exactly $21$ of the $42$ nonzero classes mod $43$). \\
% $-67$ & $-67$  & $\left(\frac{-67}{p}\right)=+1$ (exactly $33$ of the $66$ nonzero classes mod $67$). \\
% $-163$& $-163$ & $\left(\frac{-163}{p}\right)=+1$ (exactly $81$ of the $162$ nonzero classes mod $163$). \\
% \hline
% \end{tabular}
% \caption{Heegner (class number $1$) fields still have infinitely many split primes. 
% UFD means unique factorisation of \emph{elements}, not “no splitting.”}
% \end{table}
% 
% 
% \medskip
% The Gaussian $4n+1/4n+3$ story for $\mathbb{Q}(i)$ with $\Delta=-4$ is simply the most familiar member of this family—%
% the one where the congruence classes are smallest and the geometry easiest to picture—%
% but the same discriminant–residue criterion governs every quadratic field, from $-5$ and $-7$ all the way to $-163$.
% 
% 
% 
% 
% >>>>>> end of annexed block <<<<<<
\section*{Nina and John are a Reckoning}

\lettrine{A} long mahogany table stretches across the room, cluttered with half-filled notebooks, and the lumbering weight of unfinished work. The rain taps against the window as John stands arms folded ready for the sanitarium, watching the water slide down in uneven rivulets. Nina sits at the table, flipping through pages without reading a single word. John, with a small, humorless breath, \textit{"We can work separately if you’d prefer."}  

Nina, not looking up, \textit{"Would that change anything?"}  

\smallskip

John exhales sharply, the sound nearly lost in the rain. \textit{"Probably not."}  

\smallskip

\noindent
Silence stretches, old and unshaken.  

\smallskip

Nina, after many a beat, \textit{"I always wondered."}  

\smallskip

John, finally turning, \textit{"Wondered what?"}  

\smallskip

Nina, \textit{"If you realized what you were doing at the time, or was just... instinct."}  

\smallskip

John, partially unfurling \textit{"You make it sound like something deliberate."}  

\smallskip

Nina, her fingers tracing the margin of a page,  \textit{"Wasn’t it?"}  

\smallskip

\noindent
John, with something sigh-like, \textit{"You think too highly of me."}  

\smallskip

Nina, tilting her head dog-like slightly, \textit{"You hesitated."}  

\smallskip

John, tightening again, \textit{"I hesitated. And you moved forward."}  

\smallskip

Nina, now looking at him looking elsewhere,  \textit{"Someone had to."}  

\smallskip

The rain keeps falling. No answer comes, not right away.  

\smallskip

Then John, \textit{"You never had to say my name."}  

\smallskip

Nina, pressing her lips together, \textit{"Would it have made a difference?"}  
\smallskip

John, smiling devoid of warmth, \textit{"It already did."}  

\smallskip

The conversation should end there.  

\smallskip

Nina, voice now careful, \textit{"Do you think I don’t regret it?"}  

\smallskip

\noindent
John, shaking his head slightly,the weight of the past settling between them, \textit{"I think regret is cheap when the outcome has already been written."} And then after-thought like - but thoughtlessly, \textit{"Then let’s write something new."}  

\smallskip

\noindent
Her chair knowingly gnawingly scraping against the wooden floor, Nina closes the notebook in front of her.  Nina, standing, \textit{"Let’s see if you mean that this time."}  

\newpage
